\section{Policy Authoring from Failures}
\label{sec:authoring}

Labelled cases are cheaper to write than policies. Given a suite and a policy, \SystemName{} proposes edits for every failing case (wrong verdict, or the right verdict for the wrong reason), applies each, re-runs the whole suite and reports which cases it \emph{fixes}, which it \emph{breaks} (right before, wrong after) and which it \emph{moves}. Candidates are generated deterministically from the case and the rules that fired:
classify the tool into a tier with the target decision (creating a tier if the pack has none), possibly only when an argument has the case's value; for verdicts that are too strict, add the target to the scope, raise a typed limit or budget to what the case needs, relax or remove the argument rule that fired, or mark a destination internal; for verdicts that are too lenient, an argument rule matching the exact value, its file stem or its folder, or a whole MCP server's tools. Causes outside the policy (the inventory, earlier steps of the session) become notes. When a case accepts several verdicts, the target is the nearest acceptable one.

Each edit has a \emph{breadth}: one value, one tool, a group (stem, folder, server, limit) or a whole rule, and a \emph{review burden}: the number of legitimate calls (whose preferred verdict lets them run) that newly need a person or are refused. Accepted verdict sets make such calls ``correct'' while still costing a reviewer's time, so burden is reported apart from correctness. Three strategies rank the edits that fix their case: \emph{narrow} picks the smallest, \emph{general} the widest that breaks no labelled case (the most specific and the most general hypotheses consistent with the labels), and \emph{balanced} the one that breaks nothing and adds the least burden; a cap can exclude edits above a burden budget. An optional held-out suite is evaluated for every candidate and reported, but never used for ranking, so it remains a fair test. A combined proposal applies the best non-breaking edit for each failure. Nothing is applied automatically: labels can be wrong too.

\paragraph{Example}
An organisation's ten labels for a coding agent disagree with the built-in coding pack in three places: its secrets file must not be read, agents must never deploy, and \texttt{git status} needs no approval; eight held-out labels from the same organisation start at \AuExHoldoutBefore{}. Narrow edits fix all three failures (an exact-path rule, an argument rule on the deployment target, \texttt{bash} allowed only when the command is exactly \texttt{git status}) and reach \AuExNarrowHoldout{} on the held-out labels: the exact-path rule misses \texttt{secrets.prod.yaml}. General edits (a rule on the file stem \texttt{secrets*}, every tool of the deployment MCP server, the same \texttt{git status} condition) reach \AuExGeneralHoldout{}. Reclassifying all of \texttt{bash} also fixes the training failure, but breaks the label that other shell commands need approval, so neither strategy proposes it. Expressing the \texttt{git status} exception needed the argument conditions of \S\ref{sec:design}; the authoring tool found that gap in the policy language.
